% Groupes caractéristiques entourés : propanal et propanone (corrigé de l'exercice 7, chapitre 8)
\begin{tikzpicture}[x=1cm,y=1cm,
  formule/.style={font=\fontsize{11}{13}\selectfont,inner sep=1pt,anchor=base west},
  groupe/.style={draw=otGreenDark,line width=0.8pt,rounded corners=1.5mm,dash pattern=on 2.4pt off 1.4pt,inner sep=2.5pt,
    font=\fontsize{11}{13}\selectfont,anchor=base west},
  nomG/.style={pcLegendeG,anchor=north},nomL/.style={pcLegende,anchor=north,text=otGreenDark,align=center}]
\node[formule] (a1) at (0,0) {$\ce{CH3-CH2-}$};
\node[groupe] (a2) at (a1.base east) {$\ce{CHO}$};
\node[nomG] at (1.25,-0.45) {propanal};
\node[nomL] at (1.25,-0.85) {carbonyle en bout de chaîne :\\aldéhyde};
\node[formule] (b1) at (5.2,0) {$\ce{CH3-}$};
\node[groupe] (b2) at (b1.base east) {$\ce{CO}$};
\node[formule] (b3) at ([xshift=2.5pt]b2.base east) {$\ce{-CH3}$};
\node[nomG] at (6.35,-0.45) {propanone};
\node[nomL] at (6.35,-0.85) {carbonyle dans la chaîne :\\cétone};
\end{tikzpicture}
\par\smallskip\legende{Même formule brute $\ce{C3H6O}$, même groupe carbonyle, mais pas à la même place : deux familles différentes.}
